\section{Simulation Details}
\label{app:simulation-details}

This appendix summarizes the simulation design and implementation choices used
in Section~\ref{sec:experiments}. The full experiment scripts, configuration
choices, and plotting code are included in the accompanying repository.

\paragraph{Environment.}
We adapt the synthetic environment of \citet{geng2020deep}. States are
continuous with dimension \(p=5\), the action set has \(|\mathcal A|=5\)
actions, and the discount factor is \(\gamma=0.9\). We use an infinite-horizon
stationary data-generating process and generate finite trajectories of length
\(10\) for offline training and evaluation. The normalization is the anchor
constraint \(r(s,a^\dagger)=0\), implemented as anchor-action normalization
with \(a^\dagger=0\) and \(g(s)=0\). Train and test sets within each seed are
generated from the same underlying MDP parameters, with independent rollouts.

\paragraph{Behavior policy and transitions.}
The logged policy is generated from a soft \(Q\)-planner with \(\alpha=1\), so
that the DeepPQR log-policy identity is correctly matched. The planner induces
a heterogeneous but non-degenerate behavior policy, and we vary anchor support
by shifting the anchor-action logit. Transitions are deterministic conditional
on state and action up to boundary handling: actions induce action-specific
state shifts, and trajectories that leave the bounded state region are reset
uniformly within that region. We clip estimated action probabilities to
\([0.01,0.99]\) and renormalize before using them in either GenPQR or DeepPQR.

\paragraph{Experiment 1.}
The matched DeepPQR-vs.-GenPQR comparison uses shared AIRL policy estimation
and neural downstream approximation. We consider three settings:
\texttt{(200, -1.0)}, \texttt{(1000, -1.0)}, and \texttt{(2500, 0.0)}, where
the first entry is the number of training trajectories and the second is the
anchor-logit shift. These correspond to low-sample rare-anchor, mid-sample
rare-anchor, and high-sample common-anchor regimes. Each setting uses \(300\)
test trajectories and \(100\) random seeds.

\paragraph{Experiment 2.}
The higher-sample comparison fixes \(1000\) training trajectories, \(300\) test
trajectories, and anchor-logit shift \(0.0\), again over \(100\) seeds. We
compare DeepPQR, GenPQR with two policy estimators and two \(Q\)-estimators,
and several reward-recovery baselines that are standard in the IRL literature.

\paragraph{Policy estimators.}
AIRL uses a standard action-independent reward network and potential network,
both implemented as two-layer ReLU MLPs with hidden widths \((64,64)\), Adam
step size \(10^{-3}\), \(60\) behavior-cloning warm-start epochs, and \(80\)
adversarial updates in the paper experiments. Behavior cloning uses the same
\((64,64)\) MLP architecture and \(40\) epochs in Experiment~2. MaxEnt-IRL is
implemented as a neural softmax-\(Q\) model with hidden widths \((128,128)\)
and \(150\) gradient steps. These architectures were chosen to be standard for
low-dimensional synthetic control problems and to remain stable across seeds;
larger networks did not materially improve performance in pilot runs.

\paragraph{\(Q\)-evaluation.}
Neural FQE uses a dueling-style MLP \(Q(s,a)=V(s)+A(s,a)-\bar A(s)\) with
hidden widths \((128,128)\), Adam step size \(5\times 10^{-3}\), \(8\) Bellman
iterations, and \(4\) epochs of regression per iteration. Boosted FQE uses
LightGBM directly, with \(4\) outer Bellman iterations and \(30\) boosting
rounds per iteration, learning rate \(0.05\), \(32\) leaves, and minimum leaf
size \(20\). We selected these settings to balance Bellman-fit accuracy,
runtime, and seed-to-seed stability; the boosted configuration follows the
repository's adapted FQE implementation and avoids long inner re-fitting loops.

\paragraph{DeepPQR and baselines.}
DeepPQR uses the same shared AIRL policy estimate as GenPQR in the matched
comparison. It then estimates the anchor \(Q\)-function on the anchor-action
subset, reconstructs the full \(Q\)-function from log-policy ratios, and fits
the final reward-regression network. The AIRL state-reward baseline uses the
AIRL reward head directly; the log-policy pseudo-reward baseline uses
\(\log \hat\pi(a\mid s)-g(s)\); MaxEnt-IRL uses its learned \(Q\)-surrogate as
reward; and the linear reward baseline fits a strictly linear state-action
reward model with action-specific coefficients and ridge regularization.

\paragraph{Reporting.}
For each method we report reward mean squared error, reward correlation with
the ground-truth normalized reward, held-out policy negative log-likelihood
when applicable, and wall-clock runtime. Confidence intervals are normal
approximation intervals based on the \(100\) seed-level estimates.
